
\subsection*{5.3. Remarks}
\phantomsection
\label{sga3:VIIB:5.3}
\addcontentsline{toc}{subsection}{5.3. Remarks}

\noindent\SGAInlineRefTarget{sga3:VIIB:5.3:item:a}\textup{(a)}
Call a profinite \(k\)-algebra \emph{stationary} if it is the
completed tensor product of a family of truncated formal power
series algebras.\footnote{\label{sga3:VIIB:nde:165}The author
doubtless had in mind a completed tensor product
\[
  \mathcal A
  =
  \mathop{\widehat{\bigotimes}}\nolimits^{c}_{n\in\mathbb N^*}
  \frac{k[[\omega'_n]]}
       {\bigl((x^{p^n})\bigr)_{x\in\omega'_n}},
\]
where the \(\omega'_n\) are arbitrary pseudocompact vector
spaces. Then \(\omega=\omega_{\mathcal A/k}\) identifies with
\(\prod_n\omega'_n\), and the filtration is
\(\omega_n=\prod_{i=1}^{n}\omega'_i\); this is the case of
\SGARefLink{sga3:VIIB:5.2.1.B}{5.2.1.B}. For this reason these algebras would be better called
\emph{stable}, rather than \emph{stationary}; cf.\ \textup{[Di73]},
II, \(\S\)\SGARefLink{sga3:VIIB:2.9}{2.9}, p.~75. If, for example,
\(\dim\omega'_n=1\) for every \(n\in\mathbb N^*\), then
\(\mathcal A\) represents the functor that assigns to
\(C\in\operatorname{Alf}/k\) the sequences
\((x_n)_{n\in\mathbb N^*}\) of elements of \(C\) such that
\(x_n^{p^n}=0\) and \(x_n=0\) for all sufficiently large \(n\).
This case, where \(\omega=\prod_{n\in\mathbb N^*}\omega'_n\),
corresponds in the dual setting of connected cocommutative Hopf
algebras to the case studied by M.~E.~Sweedler; cf.\
\textup{[Sw67]}, Theorem~3.}

Let \(G\) be an infinitesimal formal \(k\)-group and let \(B\)
be the affine algebra of a homogeneous space of \(G\). It follows
from \SGARefLink{sga3:VIIB:5.2:theorem}{Theorem~5.2} that, for every \(r\in\mathbb N\), the algebra
\[
  B
  \big/
  \bigl((x^{p^r})\bigr)_{x\in\mathfrak r(B)}
\]
is stationary. In particular, \(B\) is an inverse limit of
stationary algebras.\footnote{\label{sga3:VIIB:nde:166}Such an
inverse limit need not itself be a stable profinite
\(k\)-algebra in the sense of the preceding editor's note. For
example, let \(S\) be the ordinary \(k\)-vector space of
sequences \((u_1,u_2,\ldots)\) of elements of \(k\), and put
\(\omega=S^*\). Then \(\omega\) is the direct sum in
\(\operatorname{PC}(k)\) of copies \(kx_n\) of \(k\), for
\(n\in\mathbb N^*\), where \(x_n(u)=u_n\). Thus this is the case
of \SGARefLink{sga3:VIIB:5.2.1.A}{5.2.1.A}. The algebra
\[
  \mathcal A
  =
  k[[\omega]]
  \big/
  \bigl((x_n^{p^n})\bigr)_{n\in\mathbb N^*}
\]
has \(\omega_{\mathcal A/k}=\omega\) and
\(\omega_n=\prod_{i=1}^{n}kx_i\), but it is not stable:
\(\operatorname{Spf}(\mathcal A)\) represents the functor that
assigns to every \(C\in\operatorname{Alf}/k\) the infinite
sequences \((x_n)_{n\in\mathbb N^*}\) of elements of \(C\) such
that \(x_n^{p^n}=0\) for every \(n\).}

\medskip
\noindent\textup{(b)}
I do not know whether, with the notation of \SGARefLink{sga3:VIIB:5.2.1:section-proposition}{5.2.1}, \(A\) and
\(B\) can be chosen so that there exists no section
\[
  \sigma:\omega\longrightarrow I
\]
compatible with the filtrations.\footnote{%
\label{sga3:VIIB:nde:167}The editors do not know either, outside
the cases considered in \SGARefLink{sga3:VIIB:5.2.1.A}{5.2.1.A} and \SGARefLink{sga3:VIIB:5.2.1.B}{5.2.1.B}.}
Nevertheless, \(\omega\) may be any pseudocompact vector space
filtered by an increasing sequence of closed subspaces. Indeed,
given
\[
  \omega_1\subset\omega_2\subset\cdots\subset\omega,
\]
one can define on
\[
  B=A
  =
  k[[\omega]]
  \big/
  \bigl((x^{p^r})\bigr)_{x\in\omega_r}
\]
a diagonal morphism
\[
  \Delta_A:A\longrightarrow A\widehat\otimes_kA
\]
satisfying conditions \textup{\SGARefLink{sga3:vii:paragraph:5-1:equation:i}{(i)}}, \textup{\SGARefLink{sga3:vii:paragraph:5-1:equation:ii}{(ii)}}, and
\textup{(iii)} of \SGARefLink{sga3:VIIB:2.1}{2.1}. It suffices to put
\[
  \Delta_A(y)
  =
  y\widehat\otimes1+1\widehat\otimes y
\]
when \(y\) is the image in \(A\) of an element of \(\omega\).

\subsection*{5.4. Corollary}
\phantomsection
\label{sga3:VIIB:5.4}
\addcontentsline{toc}{subsection}{5.4. Corollary}

\noindent\SGAInlineRefTarget{sga3:vii:corollary:5-4}\textbf{Corollary 5.4.}
{\itshape Let \(G\) be an algebraic group over a perfect field
\(k\) of characteristic \(p>0\), let \(H\) be an algebraic
subgroup of \(G\), let \(e\) be the image of the identity element
of \(G\) in \(G/H\), and let \(A\) be the local ring of \(G/H\)
at \(e\). Then \(\widehat A\) is isomorphic to an algebra of the
form
\[
  k[[X_1,\ldots,X_r,\ldots,X_s]]
  \big/
  (X_1^{p^{n_1}},\ldots,X_r^{p^{n_r}}).
\]}

\medskip
\noindent\textit{Proof.}
Consider the infinitesimal formal groups
\[
  \widehat G
  =
  \operatorname{Spf}(\widehat{\mathcal O}_{G,e}),
  \qquad
  \widehat H
  =
  \operatorname{Spf}(\widehat{\mathcal O}_{H,e}).
\]
By \SGARefLink{sga3:VIIB:1.3.4}{1.3.4}, the completion \(\widehat A\) of
\[
  A=\mathcal O_{G/H,e}
\]
is isomorphic to
\[
  \mathscr A(\widehat G/\widehat H),
\]
and the corollary follows from \SGARefLink{sga3:VIIB:5.2:theorem}{Theorem~5.2}\textup{\SGARefLink{sga3:VIIB:5.2:theorem:item:ii}{(ii)}}.\footnote{%
\label{sga3:VIIB:nde:168}See also \textup{[DG70]},
\(\S\)III.3, Theorem~6.1.}
